\documentclass[10pt,a4paper]{amsart}
\usepackage[margin=19mm]{geometry}
\usepackage{amsmath,amssymb,mathtools}
\usepackage[scr=rsfs,cal=euler]{mathalfa}
\usepackage{xfrac}
\usepackage[unicode,hidelinks,bookmarks=false]{hyperref}
\newcommand{\ie}{i.e.\ }
\newcommand{\cf}{cf.\ }
\newcommand{\resp}{respectively\ }
\newcommand{\Subsection}{\subsection}
\providecommand{\nobreakdash}{\nobreak\hbox{-}}
\setlength{\emergencystretch}{2em}
\multlinegap0pt
\usepackage{amssymb}
\usepackage{enumerate}
\newtheorem{lemma}{Lemma}
\newtheorem{theorem}{Theorem}
\title{Distinguishing adjacent vertices by ordering edges}
\author{Aleksandra Gorzkowska, Jakub Kwa\'sny}
\date{Source: arXiv:2609.11832v1 (CC BY 4.0)}
\begin{document}
\maketitle

\noindent\emph{Source and licence.} Distinguishing adjacent vertices by ordering edges. Version arXiv:2609.11832v1 (CC BY 4.0), distributed by arXiv under the Creative Commons Attribution 4.0 International licence, http://creativecommons.org/licenses/by/4.0/, as recorded in the arXiv licence metadata for this exact version and shown on the abstract page https://arxiv.org/abs/2609.11832v1. Full licence notice: \texttt{licenses/arxiv-2609.11832-CC-BY-4.0.txt}.\par\smallskip

\noindent\textit{Editorial scope.} Counted unit: the article's theorem safe-edge (source0.tex lines 132--133). The article's complete original proof of it is delivered with it (source0.tex lines 134--146, one \texttt{proof} environment of 13 lines). No mathematics is written, repaired or re-derived: the statement and the proof are the article's own lines, in the article's own notation, and no retained line is substituted or re-flowed.  Retained uncounted as the source-local context the proof needs: lines 83--84.  Nothing was omitted from any retained span, so the package carries no omission record.  No retained line contains a citation, so the wrapper carries no bibliography and no bibliography entry.  No retained line contains a figure environment, an image-inclusion command or drawing code, so no image is delivered and the package declares no asset dependency.  Editorial formatting only, in addition: an AMS \texttt{amsart} preamble that restates the article's own declarations and macros used by the retained lines, namely the environments \texttt{lemma}, \texttt{theorem} on separate counters, together with the standard packages those lines need.  The retained environments print in this document as Lemma 1, Theorem 1 in order of appearance, whereas the article prints the same items with the numbering of its own full document.  The complete native source of the article as distributed in the arXiv e-print for this exact version is bundled unchanged as \texttt{sources/210053/source0.tex}.
\begin{lemma} Let $G$ be a connected loopless multigraph with a proper edge colouring $w$, and let $v\in V(G)$. There is a total order $\prec$ on $E(G)$ such that $f_{\prec}(x)\neq f_{\prec}(y)$ for every edge $xy\in E(G-v)$. \label{lem: single_bad_vertex}
\end{lemma}

\begin{theorem} Let $G$ be a connected graph with a proper edge colouring $w$ such that there exist two adjacent vertices with different palettes. Then $w$ is sequentially orderable. \label{thm: safe-edge}
\end{theorem}

\begin{proof}
    Choose a vertex $v$, let $V_1$ consist of all vertices with palette $S(v)$, and put $V_2=V(G)\setminus V_1$. Both sets are nonempty, since some adjacent vertices have distinct palettes. Let $M$ be the set of edges between $V_1$ and $V_2$. Connectivity ensures that $M$ is nonempty, and every edge of $M$ is safe.

    Fix a total order $\prec_M$ on $M$. For each component $Q$ of $G-M$, let $M_Q$ be the set of edges of $M$ incident with $V(Q)$. We will order $E(Q)\cup M_Q$ so that adjacent vertices in $Q$ receive distinct sequences and the restriction to $M_Q$ agrees with $\prec_M$. All sequences considered in this construction include the weights on edges of $M_Q$, so they are the full sequences of vertices of $Q$ in $G$.

    Since $G$ is connected, $M_Q$ is nonempty. Let $e_Q$ be its first edge under $\prec_M$, and let $r$ be the endpoint of $e_Q$ in $Q$. Order the edges of $M_Q\setminus\{e_Q\}$ according to $\prec_M$ and reserve them as a final block. All subsequent insertions will be made before this block. If the block is empty, there is no restriction on the insertion positions.

    Root a spanning tree of $Q$ at $r$. As in Lemma \ref{lem: single_bad_vertex}, process all vertices other than $r$ before their parents. At a vertex $u\neq r$, the edge to its parent is still unordered. Insert all but one of the unordered edges of $Q$ incident with $u$ before the final block, leaving an edge $e$. There are $d_Q(u)$ positions for $e$ relative to the other incident edges of $Q$, all available before the block. These give distinct full sequences at $u$, since $w$ is proper and the incident weights from $M_Q$ form a fixed suffix. At most $d_Q(u)-1$ neighbours in $Q$ have been processed, so a position avoiding all their sequences exists. Neighbours outside $Q$ need not be considered, since the edges joining them to $u$ are safe. Later insertions leave the sequence at $u$ unchanged.

    Once all vertices other than $r$ have been processed, every edge of $Q$ is ordered. Insert $e_Q$ before the final block. Its $d_Q(r)+1$ positions relative to the edges of $Q$ incident with $r$ give distinct full sequences, while only $d_Q(r)$ neighbour sequences in $Q$ must be avoided. Thus a suitable position exists. This also covers the case when $Q$ is a single vertex. The resulting order distinguishes all adjacent vertices in $Q$ and agrees with $\prec_M$ on $M_Q$, because $e_Q$ precedes the final block.

    Finally, the orders constructed for the components admit a common extension together with $\prec_M$. Take such a common extension to $E(G)$. Every vertex of $Q$ has all its incident edges in $E(Q)\cup M_Q$, so its sequence is preserved. All edges outside $M$ therefore have distinct endpoint sequences, and the same holds for the safe edges of $M$.
\end{proof}

\end{document}
